\documentclass[runningheads]{llncs}
\usepackage{amsmath,amssymb,booktabs,graphicx}
\title{Multi-Representation Trajectory Learning for Vietnamese Air-Writing Recognition}
\author{Anonymous Authors}
\institute{Anonymous Institution}
\begin{document}
\maketitle

\begin{abstract}
Air-writing recognition must preserve both the geometry of a trajectory and
the motion that produced it. We present a trajectory-aware character
recognizer for Vietnamese that combines spatial coordinates, motion dynamics,
and multi-scale Fourier features through an adaptive gated fusion module. A
Conformer encoder with learned relative-position bias and a CTC objective then
maps the fused sequence to Vietnamese character strings. We construct a
source-disjoint evaluation protocol for 22,754 trajectories covering 660
labels and report character error rate, word error rate, exact sequence
accuracy, and Vietnamese-specific error statistics. The final numerical
results are generated by the reproducible runs in this repository and will be
inserted from `artifacts/experiments/summary.csv`.
\keywords{air-writing \and trajectory representation \and Conformer \and CTC \and Vietnamese}
\end{abstract}

\section{Introduction}
Air-writing provides a natural interface for text entry without a physical
surface, but camera-tracked trajectories vary in position, scale, orientation,
and writing speed. Character-level Vietnamese recognition adds a dense set of
accented vowels and consonants whose short strokes are easy to delete.
Existing sequence recognizers commonly use coordinates or motion descriptors
alone. Coordinates retain global structure but are sensitive to spatial
variation, whereas derivatives emphasize motion while discarding absolute
shape. This work asks how complementary trajectory representations should be
combined for Vietnamese air-writing.

Our contributions are: (i) a joint spatial, dynamic, and multi-scale Fourier
trajectory representation; (ii) adaptive per-time-step gated fusion followed
by a relative-position Conformer-CTC recognizer; and (iii) a leakage-safe
evaluation with architecture, representation, fusion, preprocessing,
augmentation, robustness, and Vietnamese error analyses.

\section{Method}
Given a trajectory $P=\{(x_t,y_t)\}_{t=1}^T$, we smooth, interpolate,
centroid-center, globally scale, and resample it to $T_r$ points. The spatial
branch is $S_t=[x_t,y_t]$. The dynamic branch contains displacement, speed,
unit direction, acceleration, and signed curvature:
\begin{equation}
D_t=[\Delta x_t,\Delta y_t,v_t,d^x_t,d^y_t,a_t,\kappa_t].
\end{equation}
The spectral branch uses $K$ dyadic frequencies,
$F_t=[\sin(2^k\pi p_t),\cos(2^k\pi p_t)]_{k=0}^{K-1}$.
Each branch is embedded independently to $d$ dimensions. For gated fusion,
$\alpha_t=\operatorname{softmax}(W[h^S_t;h^D_t;h^F_t])$ and
$h_t=\sum_b\alpha^b_t h^b_t$. The fused sequence is processed by Conformer
blocks. Relative self-attention adds a learned bias indexed by clipped
relative distance; the output projection is trained with CTC.

\section{Experimental Protocol}
The VNI air-writing collection contains 22,760 CSV trajectories in 660 label
directories. We remove six exact duplicate files and freeze the manifest
before training. Repeated numeric source IDs are partitioned globally into
40/5/5 train/validation/test IDs, giving 18,202/2,276/2,276 samples and every
label in each split. The data contains no writer metadata, so writer
independence is not claimed. Augmentation is applied only to training data.
The primary metric is character error rate; word error rate, exact sequence
accuracy, insertion/deletion/substitution counts, length buckets, and
per-character Vietnamese errors are also reported.

\section{Results}
Tables and plots are generated from `summary.csv`, prediction JSONL files,
and `robustness.json`. No result is reported from the earlier protocol that
augmented validation or test samples. The main comparison includes BiLSTM,
GRU, TCN, Transformer, Conformer, and the proposed adaptive fusion model;
ablations isolate spatial, dynamic, Fourier, positional, augmentation, and
resampling choices. In the completed seed-42 Fourier-scale sweep, $K=2$ is
selected on validation and obtains 1.77\% CER, 5.55\% WER, and 94.77\%
exact-sequence accuracy on the held-out test split. The default $K=4$ gated
configuration obtains 2.36\% CER, while the coordinate-plus-dynamics
Conformer baseline obtains 1.85\% CER; the absolute-position Transformer
baseline obtains 11.79\% CER. These values are included as an
initial table; the remaining multi-seed and preprocessing runs are retained
in the Modal results volume and must be merged before the final submission.

\section{Conclusion}
This study evaluates trajectory representation as the central design choice
for Vietnamese air-writing recognition. The released manifest, training code,
checkpoints, and predictions make the final claims auditable and reproducible.

\end{document}
